\section{Local divisibility and finite row sums}\label{app:local}\label{supp:local}
The local estimates have two separate ingredients. A change of polynomial
basis makes divisibility visible in each row. A counting argument then bounds
how much of that divisibility can be lost to the derivative orders in a
selected minor. We give both arguments here, including the finite corrections
needed for the small primes.

Throughout this appendix, $K,S,J,n$ are positive integers, $R=KS$, and
\[
 \Lambda=\{(r,s)\in\NN^2:r+Ks<R\},\qquad
 M=\frac{R(S+1)}2,\qquad W=\frac{R^2S}3,\qquad
 \theta=\frac{M}{nJ}\leq1.
\]
The nodes $z_1,\ldots,z_n$ may be any elements of $\OO_D$. Columns are indexed
by $(i,j)$ with $1\leq i\leq n$ and $0\leq j<J$, and $C$ denotes any selection
of $M$ columns. Thus each derivative order occurs in $C$ at most $n$ times.
The determinant $\Delta_C$ is formed with the Hasse derivatives in
\eqref{eq:jet-matrix}; differentiation is in $z$, with $w$ held fixed.

\subsection{Passing from entries to a determinant}
The following elementary observation allows different derivative estimates
to be used on different sets of rows.

\begin{lemma}\label{lem:local-determinant}
Let $A=(a_{rc})$ be a square matrix over a discretely valued field, and let
$v$ be its valuation. Suppose that $E$ is a set of rows and
\[
 v(a_{rc})\geq a_r-b_c-\boldsymbol1_{r\in E}d_c,
 \qquad d_c\geq0.
\]
Then
\[
 v(\det A)\geq\sum_r a_r-\sum_c b_c
 -\max_{\substack{I\subseteq C\\|I|=|E|}}\sum_{c\in I}d_c.
\]
An integral unitriangular change of rows does not change the determinant.
If columns of order $j$ have cost $d(j)$, each order is available at most
$n$ times, and $d(j)$ is nondecreasing, the maximum cost of $m$ columns is
bounded by filling the orders $J-1,J-2,\ldots$ in that order, with capacity
$n$ at each order. The minimum of a nondecreasing cost is bounded from below
by filling orders $0,1,\ldots$ in the same way.
\end{lemma}
\begin{proof}
Each term of the determinant uses every row and every column once. The
columns paired with rows in $E$ form a set of size $|E|$, so the displayed
entry estimate gives the asserted lower bound for every determinant term.
The valuation of a sum is at least the minimum valuation of its terms.
A unitriangular row transformation has determinant one. Finally, if a
selected column has smaller cost than an available column, exchanging them
cannot decrease the total cost. Repeating this exchange proves the capacity
assertions, and reversing the exchanges proves the assertion about minima.
\end{proof}

\subsection{Two finite row sums}\label{supp:rows}
For integers $N\geq0$ and $m\geq1$, define
\[
 T(N,m)=\sum_{s=0}^{N-1}(N-s)\floor{s/m}.
\]
An empty sum is zero. Expressing $\floor{s/m}$ as the number of positive
multiples of $m$ not exceeding $s$, and interchanging the two finite sums,
gives
\begin{equation}\label{eq:T-floor}
 T(N,m)=\frac12\sum_{\substack{a\geq1\\am<N}}(N-am)(N-am+1).
\end{equation}
There are $K(S-s)$ rows with a given second coordinate $s$. Consequently
\[
 \sum_{(r,s)\in\Lambda}\sum_{a\geq1}\floor{s/q^{2a}}
 =K\sum_{a\geq1}T(S,q^{2a}).
\]
In particular, this sum is finite even when written with no upper limit.

The corresponding sum in the first coordinate is conveniently bounded by
\begin{equation}\label{eq:Sr}
 S_r(R,K,x)=\sum_{a=1}^{\floor{R/x}}(R-ax)(R+K-ax),\qquad x>0.
\end{equation}
For an integer $q\geq1$, put $Q=\floor{R/q}$ and $c_a=aq\bmod K$.
Then the exact identity is
\begin{equation}\label{eq:local-r-sum}
 2K\sum_{\Lambda}\floor{r/q}
 =S_r(R,K,q)+\sum_{a=1}^{Q}c_a(K-c_a)
 \geq S_r(R,K,q).
\end{equation}
To see this, count the rows with $r\geq aq$. Write $aq=Kb+c$, where
$0\leq c<K$, and put $u=S-b$. For each $s=0,\ldots,u-1$ there are
$K(u-s)-c$ such rows. Their number is $Ku(u+1)/2-cu$.
Multiplication by $2K$ gives
$(R-aq)(R+K-aq)+c(K-c)$. Summing over $a$ proves
\eqref{eq:local-r-sum}.

We will also need the shape of this weight as its argument varies. On an
interval on which $Q=\floor{R/x}$ is fixed, summing $a$ and $a^2$ gives
\begin{equation}\label{eq:local-Sr-polynomial}
 S_r(R,K,x)=QR(R+K)
 -\frac{Q(Q+1)(2R+K)}2x
 +\frac{Q(Q+1)(2Q+1)}6x^2.
\end{equation}
At an endpoint $x=R/Q$, the summand that disappears is zero, so $S_r$ is
continuous. If $u=R-Qx$, differentiation of this polynomial yields
\[
 -\frac{d}{dx}\bigl(xS_r(R,K,x)\bigr)
 =Q\left(\frac{Q-1}{2}x^2+u(2x-u)+K(x-u)\right).
\]
For $0<x<R$ we have $Q\geq1$ and $0\leq u<x$, and the right side is
positive. Since $\log x/x$ is positive and decreasing for $x\geq3$, it
follows that
\begin{equation}\label{eq:local-weight-decreasing}
 x\longmapsto\frac{S_r(R,K,x)\log x}{KW}
 \quad\hbox{is decreasing for }3\leq x<R.
\end{equation}
It is zero for $x\geq R$. This is the monotonicity used when replacing a
possible prime norm by a lower endpoint for that norm.

For the unbounded range, an explicit lower bound is more useful than a
piecewise polynomial. Set $\vartheta=\{R/q\}$ and substitute
$Q=R/q-\vartheta$ into \eqref{eq:local-Sr-polynomial}. Expanding gives
\begin{align*}
 S_r(R,K,q)={}&\frac{R^3}{3q}-\frac{R^2}2
 +\frac{KR^2}{2q}+\frac{qR}6-\frac{KR}2\\
 &-\frac{q^2}{6}\vartheta(1-\vartheta)(1-2\vartheta)
 +\frac{Kq}{2}\vartheta(1-\vartheta).
\end{align*}
The last term is nonnegative, and
$|\vartheta(1-\vartheta)(1-2\vartheta)|\leq1/(6\sqrt3)$.
Using $KW=R^3/3$, we obtain
\begin{equation}\label{eq:local-Bernoulli}
 \frac{S_r(R,K,q)}{KW}\geq
 \frac1q-\frac{3}{2R}+\frac{3K}{2Rq}
 +\frac{q}{2R^2}-\frac{3K}{2R^2}
 -\frac{q^2}{12\sqrt3 R^3}.
\end{equation}
The cubic expression in $\vartheta$ is the usual third Bernoulli-polynomial
remainder; the displayed calculation fixes both its sign and its constant.

\subsection{The basis for a small inert prime}
Fix an inert prime $q<K$, and write
\[
 h=h_q(r,s)=\min\left(\floor{r/q},
       \floor{(R-1-r-Ks)/(K-q)}\right),\qquad
 \tau_q(s)=\sum_{a\geq1}\floor{s/q^{2a}}.
\]
The first bound on $h$ permits $h$ copies of $z^q$ to be replaced by $w$.
The second ensures that the resulting monomials stay inside $\Lambda$.
For fixed $s$ put $t=S-s$. The inequalities $h\geq a$ say exactly
$qa\leq r<Kt-(K-q)a$, so there are $K(t-a)$ such rows when $a<t$,
and none otherwise. Summing these superlevel counts gives
\begin{equation}\label{eq:local-h-sum}
 2\sum_\Lambda h_q(r,s)=\frac{R(S^2-1)}3=W_2.
\end{equation}

We first describe the rows on which the two derivative estimates will be
used. In the part $0\leq r<q(S-s)$, one has $h=\floor{r/q}$, so
$0\leq r-qh<q$. In the remaining part, write
\[
 R-1-r-Ks=(K-q)h+e,\qquad 0\leq e<K-q.
\]
Here $h$ is the second member of the minimum. If
$d=S-s-1-h$ and $c=K-1-e$, then
\begin{equation}\label{eq:local-strip}
 r-qh=Kd+c,\qquad 0\leq d<S,\quad q\leq c<K.
\end{equation}
Conversely, for any such $d,c$, the values
\[
 0\leq s<S-d,\qquad h=S-d-1-s,\qquad r=qh+Kd+c
\]
give exactly these rows. Thus for each pair $(d,c)$ there is a strip of
length $t=S-d$, and there are $K-q$ strips of every length
$1,\ldots,S$. The number
\[
 m=\floor{(r-qh)/q^2}=\floor{(Kd+c)/q^2}
\]
is constant on a strip. This constancy is what makes the finite estimate
below possible.

Choose an integer $L$ with $0\leq L\leq S$, and let $E_L$ consist of the
strips of length at most $L$. In particular,
\begin{equation}\label{eq:local-cut-size}
 m_L:=|E_L|=\frac{(K-q)L(L+1)}2,\qquad \varpi_L:=\frac{m_L}{M}.
\end{equation}
On these rows we
will use all available factors $w^q-z$. On the other rows we keep the tower
in $w$, which has no derivative cost.

Let $T_a$ be the inert tower from \eqref{eq:tower-value}. Writing
$s=\sum_{a\geq0}d_aq^{2a}$, put
\[
 B_s(w)=\prod_{a\geq1}T_a(w)^{d_a},\qquad
 u(s)=\floor{s/q}\bmod q,\qquad r_0=r-qh-q^2m.
\]
Thus $B_s$ is monic of degree $q^2\floor{s/q^2}$ and
$v_q(B_s(\bar z_i))\geq\tau_q(s)$. Define
\begin{equation}\label{eq:local-mixed-basis}
 P_{r,s}(z,w)=z^{r_0}(z^{q^2}-z)^m(z^q-w)^h w^{s\bmod q}
 \begin{cases}
 (w^q-z)^{\floor{s/q}},&(r,s)\in E_L,\\
 (w^q-z)^{u(s)}B_s(w),&(r,s)\notin E_L.
 \end{cases}
\end{equation}
In both cases the leading choice from the factors is $z^rw^s$, with
coefficient one.

Here are the two checks needed to use these polynomials as row operations.
First, replacing $z^q$ by $w$ increases the weight $r+Ks$ by $K-q$, but
there are at most $h$ such replacements. Their total increase is at most
$R-1-r-Ks$. Every other replacement in
\eqref{eq:local-mixed-basis} decreases the weight. All monomials therefore
remain in $\Lambda$. Second, order monomials lexicographically by
$(qr+s,s)$. Replacing $z^{q^2}$ by $z$, or $z^q$ by $w$, strictly decreases
the first coordinate. Replacing $w^q$ by $z$ preserves the first coordinate
and decreases the second. A lower term of $B_s$ decreases the first
coordinate. Hence the change from $z^rw^s$ to $P_{r,s}$ is integral and
unitriangular. It preserves every selected minor.

\subsection{Valuations and derivative costs}
At a node, Frobenius and conjugation agree modulo an inert prime. Taylor
expansion in $z$ consequently has the form
\[
 (z_i+t)^q-\bar z_i=qU(t)+t^q,\qquad
 \bar z_i^q-(z_i+t)=qV-t,
\]
\[
 (z_i+t)^{q^2}-(z_i+t)=qH(t)-t+t^{q^2},
\]
with $U,H\in\OO_D[t]$ and $V\in\OO_D$. Let $a$ denote the combined number
of factors of the latter two types: $a=m+\floor{s/q}$ on $E_L$, and
$a=m+u(s)$ off $E_L$.

To lose a factor of $q$ from one of the $h$ factors of the first type, a
Taylor term must spend at least $q$ degrees in $t$. To lose one from any of
the other $a$ factors, it must spend at least one degree. If $a_1$ and $a_2$
are the respective numbers of such choices in a term of degree $j$, then
$qa_1+a_2\leq j$ and $a_2\leq a$. Its valuation is at least
$h+a-a_1-a_2$. These two inequalities imply, respectively,
\[
 h+a-a_1-a_2\geq h+a-j,
 \qquad
 h+a-a_1-a_2\geq h-\floor{(j-a)/q}.
\]
All factors not yet counted have integral coefficients. Thus the entry
valuation, before adding the tower contribution, is at least
\begin{equation}\label{eq:local-entry}
 \max\left(0,h+a-j,h-\floor{(j-a)/q}\right).
\end{equation}

Set $\delta_q(j)=j-\floor{j/q}$. For a row outside $E_L$, the last member
of \eqref{eq:local-entry}, together with the tower, gives
\begin{equation}\label{eq:local-entry-low}
 v_q\bigl(\partial_z^{[j]}P_{r,s}(z_i,\bar z_i)\bigr)
 \geq h+\tau_q(s)+\floor{(m+u(s))/q}-\floor{j/q}.
\end{equation}
For a row in $E_L$, the middle member gives instead
\begin{align}
 v_q\bigl(\partial_z^{[j]}P_{r,s}(z_i,\bar z_i)\bigr)
 \geq{}&h+\tau_q(s)
       +\bigl(m+\floor{s/q}-\tau_q(s)\bigr)\notag\\
       &-\floor{j/q}-\delta_q(j).
 \label{eq:local-entry-high}
\end{align}
Both estimates have the same baseline row contribution $h+\tau_q(s)$
and the same column cost $\floor{j/q}$. The extra row contribution is
nonnegative in both cases, since $\tau_q(s)\leq\floor{s/q}$.
Only rows in $E_L$ incur the additional cost $\delta_q(j)$.

\subsection{Summing the extra row contribution}
For the rows outside the cut, the following finite form of the floor-sum
identity controls an incomplete strip.

\begin{lemma}\label{lem:local-Hermite}
For integers $m,t\geq0$ and $q\geq2$,
\[
 \sum_{s=0}^{t-1}\floor{\bigl(m+(\floor{s/q}\bmod q)\bigr)/q}
 \geq\frac{tm}{q}-\frac{q^2}{4}.
\]
\end{lemma}
\begin{proof}
Write $m=aq+b$, with $0\leq b<q$. On each complete block of $q^2$
consecutive values of $s$, the summand is $a$ for $q(q-b)$ values and
$a+1$ for $qb$ values. Its average is exactly $m/q$.
It remains to examine a final block of length $kq+v$, with
$0\leq k,v<q$. If $k<q-b$, its deficit from this average is
$b(k+v/q)\leq b(q-b)$. If $k\geq q-b$, the deficit is
$(q-b)(q-k-v/q)\leq b(q-b)$. Both bounds are at most $q^2/4$.
The case $b=0$ has zero deficit.
\end{proof}

We now sum over the strips in \eqref{eq:local-strip}. Put
\[
 c_{K,q}=\frac{K+q-1}{2}-(q^2-1).
\]
For a fixed length $t$, summing
$\floor{(K(S-t)+c)/q^2}\geq(K(S-t)+c-(q^2-1))/q^2$
over $c=q,\ldots,K-1$ gives
\begin{equation}\label{eq:local-strip-mass}
 \sum_{c=q}^{K-1}\floor{(K(S-t)+c)/q^2}
 \geq\frac{K-q}{q^2}\bigl(K(S-t)+c_{K,q}\bigr).
\end{equation}
The right side is allowed to be negative; the exact extra contribution is
still nonnegative.

Define
\begin{align*}
 A_q(L)&=\frac{2(K-q)}{q^2}\sum_{t=1}^{L}
              t\bigl(K(S-t)+c_{K,q}\bigr),\\
 B_q(L)&=\frac{2(K-q)}{q^3}\sum_{t=L+1}^{S}
              t\bigl(K(S-t)+c_{K,q}\bigr),\\
 E_q(L)&=A_q(L)+B_q(L)-\frac{q^2}{2}(S-L)(K-q)\\
 &\hspace{12mm}+2(K-q)\left(T(L,q)-\sum_{a\geq1}T(L,q^{2a})\right).
\end{align*}
Twice the sum of the extra row contributions in
\eqref{eq:local-entry-low}--\eqref{eq:local-entry-high} is at least $E_q(L)$.
Indeed, on each strip of length $t\leq L$, the contribution of $m$ is $tm$;
\eqref{eq:local-strip-mass} sums these contributions to at least $A_q(L)/2$.
The remaining contribution on those strips is
\[
 (K-q)\sum_{t=1}^{L}\sum_{s=0}^{t-1}
       \left(\floor{s/q}-\tau_q(s)\right)
 =(K-q)\left(T(L,q)-\sum_{a\geq1}T(L,q^{2a})\right).
\]
Here a given $s$ occurs $L-s$ times, explaining the function $T$.

For a strip of length $t>L$, Lemma~\ref{lem:local-Hermite} bounds the extra
contribution from below by $tm/q-q^2/4$. Summing the first term with
\eqref{eq:local-strip-mass} gives $B_q(L)/2$; there are
$(S-L)(K-q)$ such strips. The rows in the remaining left part of the
triangle have nonnegative extra contribution and may be omitted. These
observations account for every term of $E_q(L)$. Since the exact total
extra contribution is itself nonnegative, we have the stronger lower bound
\begin{equation}\label{eq:local-extra-positive}
 2\sum_\Lambda(\text{extra row contribution})\geq\pos{E_q(L)}.
\end{equation}
The positive part is taken before subtracting the extra derivative cost.

\subsection{The extra column cost and the retained weight}
The cost $\delta_q(j)$ is nondecreasing. Its prefix sum has a simple exact
form:
\begin{equation}\label{eq:small-column-prefix}
 \Sigma_q(B):=\sum_{j=0}^{B-1}\delta_q(j)
 =\frac{q-1}{2q}B^2-\frac{r_B(q-r_B)}{2q},
 \qquad r_B=B\bmod q.
\end{equation}
This follows by writing $B=aq+r_B$ and summing the complete blocks of the
floor function. It also holds for $B=0$.

Write $m_L=nt+r$, where $0\leq r<n$, and put $A=J-t$. By
Lemma~\ref{lem:local-determinant}, the extra derivative cost is at most
\[
 n\bigl(\Sigma_q(J)-\Sigma_q(A)\bigr)+r\delta_q(A-1).
\]
The last term is omitted when $r=0$. In particular, $A-1$ is only used
when $A\geq1$. To estimate the remainder in this formula, rewrite it as
\[
 n\Sigma_q(J)-(n-r)\Sigma_q(A)-r\Sigma_q(A-1).
\]
For $r=0$ the term with $A-1$ is again omitted. The remainder in $\Sigma_q(J)$
is nonpositive, whereas each of the other two remainders is bounded above
in absolute value by $\floor{q^2/4}/(2q)$. The quadratic parts give
\[
 \frac{q-1}{q}\left(m_LJ-\frac{m_L^2}{2n}\right)
 -\frac{q-1}{2q}r\left(1-\frac rn\right).
\]
Dropping the final nonpositive term proves the upper bound
\begin{equation}\label{eq:local-extra-cost}
 \frac{q-1}{q}\left(m_LJ-\frac{m_L^2}{2n}\right)
 +\boldsymbol1_{m_L>0}\frac{n\floor{q^2/4}}{2q}.
\end{equation}
If $m_L=0$, the exact cost is zero, so no remainder is charged.

Normalize twice this cost by $W$, and recall $m_L=\varpi_LM$ and
$M/(nJ)=\theta$. Since $2M/W=3(S+1)/(RS)$, the resulting expression is
\[
 C_q(L)=3\left(1+\frac1S\right)\varpi_L
       \left(1-\frac{\varpi_L\theta}{2}\right)\frac{q-1}{q}\frac JR
       +\boldsymbol1_{L>0}\frac{n\floor{q^2/4}}{qW}.
\]
We have now obtained every term of the small-prime formula.

\begin{proposition}\label{prop:local-small-prime}
For an inert prime $q<K$, define
\begin{equation}\label{eq:wkeep}
 w_{\rm keep}(q)=
 \frac{W_2+2K\sum_{a\geq1}T(S,q^{2a})}{W}
 +\max_{0\leq L\leq S}
       \left(\frac{\pos{E_q(L)}}W-C_q(L)\right).
\end{equation}
Then every selected minor satisfies
\[
 v_q(\Delta_C)\geq\frac W2w_{\rm keep}(q)
                       -\sum_{(i,j)\in C}\floor{j/q}.
\]
The maximum is over integers $L$. Using any one admissible value of $L$
in place of the maximum also gives a valid lower bound.
\end{proposition}
\begin{proof}
For the basis associated with $L$, the sum of the baseline row contributions
is, by \eqref{eq:local-h-sum} and \eqref{eq:T-floor},
\[
 \sum_\Lambda(h+\tau_q(s))
 =\frac{W_2}{2}+K\sum_{a\geq1}T(S,q^{2a}).
\]
Equation~\eqref{eq:local-extra-positive} bounds the extra row contribution
from below, and \eqref{eq:local-extra-cost} bounds its extra column cost
from above. Apply Lemma~\ref{lem:local-determinant} to
\eqref{eq:local-entry-low} and \eqref{eq:local-entry-high}. Every choice
of $L$ gives a bound for the same determinant, because all the bases are
unitriangular. Taking the best of these bounds proves the proposition.
\end{proof}

In particular, the argument includes $q=2$ when $2$ is inert in the maximal
order. The factors $1+1/S$ and the last remainder in $C_q(L)$ are exact
finite-size corrections.

\subsection{Split and ramified primes}
For completeness we give the local contributions at the other primes.
They are needed when primes beyond the guaranteed inert range are retained.
At a split prime, a tower in $w$ gives a contribution with no derivative
cost. At a ramified prime, the inert estimates remain valid in the prime
ideal valuation, but this has half the norm contribution of the inert
valuation.

Suppose first that $q$ splits. Then
$\OO_D/q\OO_D\simeq\FF_q\times\FF_q$, so $b^q-b$ is divisible by $q$
for every $b\in\OO_D$. Define monic polynomials
\[
 G_1(w)=w^q-w,\qquad
 G_{a+1}(w)=G_a(w)^q-q^{q^a-1}G_a(w),\qquad
 \mu_a=\frac{q^a-1}{q-1}.
\]
They have degrees $q^a$, and $G_a(\bar z_i)\in q^{\mu_a}\OO_D$.
Indeed, if $G_a(\bar z_i)=q^{\mu_a}b$, the next value is
$q^{q\mu_a}(b^q-b)$, divisible by $q^{q\mu_a+1}=q^{\mu_{a+1}}$.
For the base-$q$ expansion $s=\sum_{a\geq0}d_aq^a$, the polynomial
\[
 \widetilde A_s(w)=w^{d_0}\prod_{a\geq1}G_a(w)^{d_a}
\]
is monic of degree $s$ and has values divisible by
$q^{\sum_{a\geq1}\floor{s/q^a}}$. The basis
$z^r\widetilde A_s(w)$ is unitriangular on $\Lambda$, and differentiating
in $z$ does not alter this divisibility. Hence
\begin{equation}\label{eq:local-split}
 \Delta_C\in q^{K\sum_{a\geq1}T(S,q^a)}\OO_D.
\end{equation}
Its contribution to $\log|\Delta_C|^2/W$ is at least
\[
 \bigl(e(q)+s(q)\bigr)\log q,\qquad
 e(q)=\frac{2K}{W}\sum_{a\geq1}T(S,q^{2a}),\quad
 s(q)=\frac{2K}{W}\sum_{a\geq1}T(S,q^{2a-1}).
\]
There is no first-power derivative debit in this estimate. When the
arithmetic lower bound is organized with a common debit
$\sum_C\floor{j/q}$ for all primes, a split prime therefore permits that
debit to be refunded. Its minimum over $C$ is given by the capacity
principle in Lemma~\ref{lem:local-determinant}.

Now suppose that $q\geq K$ and $(q)=\mathfrak q^2$ is ramified. In our
application the ramified prime is $q=D>K$. Normalize $v_{\mathfrak q}$
by $v_{\mathfrak q}(\mathfrak q)=1$. Conjugation acts trivially on
$\OO_D/\mathfrak q\simeq\FF_q$, so $z_i^q-\bar z_i$ and
$\bar z_i^q-\bar z_i$ belong to $\mathfrak q$.
A substitute for the inert tower is
\[
 \widehat B_s(w)=(w^q-w)^{q\floor{s/q^2}}.
\]
It is monic of degree $q^2\floor{s/q^2}$, and
\[
 v_{\mathfrak q}(\widehat B_s(\bar z_i))
 \geq q\floor{s/q^2}\geq\tau_q(s).
\]
For the last inequality, put $a=\floor{s/q^2}$; then
$\tau_q(s)=a+\floor{a/q^2}+\floor{a/q^4}+\cdots
\leq aq^2/(q^2-1)\leq aq$.

Use the unitriangular basis
\[
 z^{r\bmod q}(z^q-w)^{\floor{r/q}}
 w^{s\bmod q^2}\widehat B_s(w).
\]
As in \eqref{eq:tower-valuation}, its entries have valuation at least
$\floor{r/q}+\tau_q(s)-\floor{j/q}$. Therefore
\begin{equation}\label{eq:local-ramified}
 v_{\mathfrak q}(\Delta_C)\geq
 \frac{S_r(R,K,q)}{2K}+K\sum_{a\geq1}T(S,q^{2a})
 -\sum_{(i,j)\in C}\floor{j/q}.
\end{equation}
One may replace a negative right side by zero, since the entries are
integral. Finally, an inert prime ideal has norm $q^2$, whereas a ramified
prime ideal has norm $q$. The contribution of
$v_{\mathfrak q}(\Delta_C)$ to $\log|\Delta_C|^2$ is consequently
$v_{\mathfrak q}(\Delta_C)\log q$. This explains the factor of one half
in the normalized ramified contribution used in the arithmetic comparison.
